
\documentclass{article}
\usepackage{colortbl}
\usepackage{makecell}
\usepackage{multirow}
\usepackage{supertabular}

\begin{document}

\newcounter{utterance}

\centering \large Interaction Transcript for game `codenames', experiment `opponent\_none', episode 2 with gemma{-}4{-}26b{-}a4b{-}it.
\vspace{24pt}

{ \footnotesize  \setcounter{utterance}{1}
\setlength{\tabcolsep}{0pt}
\begin{supertabular}{c@{$\;$}|p{.15\linewidth}@{}p{.15\linewidth}p{.15\linewidth}p{.15\linewidth}p{.15\linewidth}p{.15\linewidth}}
    \# & $\;$A & \multicolumn{4}{c}{Game Master} & $\;\:$B\\
    \hline

    \theutterance \stepcounter{utterance}  
    & & & \multicolumn{2}{p{0.3\linewidth}}{
        \cellcolor[rgb]{0.9,0.9,0.9}{
            \makecell[{{p{\linewidth}}}]{
                \texttt{\tiny{[GM$|$GM]}}
                {'hidden': {'team': ['barco', 'tubo', 'corbata', 'espalda', 'buceador', 'célula', 'columpio', 'lugar', 'herradura'], 'innocent': ['interruptor', 'tono', 'china', 'maestro', 'tokio', 'diente', 'dragón'], 'opponent': ['puerto', 'soldado', 'personal', 'gorra', 'gancho', 'enterrador', 'capa', 'escuela'], 'assassin': ['científico']}, 'revealed': {'team': {'team': [], 'innocent': [], 'opponent': [], 'assassin': []}, 'opponent': {'team': [], 'innocent': [], 'opponent': [], 'assassin': []}}}
            }
        }
    }
    & & \\ \\

    \theutterance \stepcounter{utterance}  
    & & \multicolumn{4}{p{0.6\linewidth}}{
        \cellcolor[rgb]{0.9,0.9,0.9}{
            \makecell[{{p{\linewidth}}}]{
                \texttt{\tiny{[P1$\langle$GM]}}
                \texttt{Proporciona una pista de una sola palabra que esté relacionada con algunas de las palabras de tu equipo de la lista a continuación. Puedes elegir con cuántas y con qué palabras se relaciona tu pista.} \\
\texttt{El objetivo es que tu compañero adivine todas las palabras de tu equipo basándose en las pistas que le des, antes de que el otro equipo haya descubierto todas las suyas. Si tu compañero adivina una palabra 'oponente' o 'distractor', su turno de adivinar termina y el equipo contrario puede hacer un turno. Si tu compañero adivina una palabra 'asesina', el juego termina y tu equipo pierde inmediatamente.} \\
\texttt{Cuantas más palabras relacione tu pista, más rápido podrás ganar el juego, pero más difíciles de adivinar serán las asociaciones para tu compañero, así que elige bien tu palabra pista. La palabra pista tiene que estar semánticamente relacionada con las palabras objetivo, no puede ser una de las palabras de las listas ni contener partes de ellas.} \\
\\ 
\texttt{Siempre da tu pista de una sola palabra y tu lista separada por comas de palabras objetivo relacionadas en el siguiente formato y haz tus respuestas lo más breves posible, nunca incluyas más texto que el requerido en este formato:} \\
\\ 
\texttt{PISTA: <PALABRA>} \\
\texttt{OBJETIVOS: <OBJETIVOS>} \\
\\ 
\texttt{Las palabras de tu equipo son: barco, tubo, corbata, espalda, buceador, célula, columpio, lugar, herradura.} \\
\texttt{Las palabras del equipo de tu oponente son: puerto, soldado, personal, gorra, gancho, enterrador, capa, escuela.} \\
\texttt{Las palabras distractor son: interruptor, tono, china, maestro, tokio, diente, dragón.} \\
\texttt{Las palabras asesinas son: científico.} \\
            }
        }
    }
    & \\ \\

    \theutterance \stepcounter{utterance}  
    & \multicolumn{4}{p{0.6\linewidth}}{
        \cellcolor[rgb]{0.9,0.9,0.9}{
            \makecell[{{p{\linewidth}}}]{
                \texttt{\tiny{[P1$\rangle$GM]}}
                \texttt{PISTA: MAR} \\
\texttt{OBJETIVOS: barco, buceador} \\
            }
        }
    }
    & & \\ \\

    \theutterance \stepcounter{utterance}  
    & & & \multicolumn{2}{p{0.3\linewidth}}{
        \cellcolor[rgb]{0.9,0.9,0.9}{
            \makecell[{{p{\linewidth}}}]{
                \texttt{\tiny{[GM$|$GM]}}
                \texttt{mar} \\
            }
        }
    }
    & & \\ \\

    \theutterance \stepcounter{utterance}  
    & & & \multicolumn{2}{p{0.3\linewidth}}{
        \cellcolor[rgb]{0.9,0.9,0.9}{
            \makecell[{{p{\linewidth}}}]{
                \texttt{\tiny{[GM$|$GM]}}
                ['barco', 'buceador']
            }
        }
    }
    & & \\ \\

    \theutterance \stepcounter{utterance}  
    & & & \multicolumn{2}{p{0.3\linewidth}}{
        \cellcolor[rgb]{0.9,0.9,0.9}{
            \makecell[{{p{\linewidth}}}]{
                \texttt{\tiny{[GM$|$GM]}}
                {'word': 'barco', 'assignment': 'team'}
            }
        }
    }
    & & \\ \\

    \theutterance \stepcounter{utterance}  
    & & & \multicolumn{2}{p{0.3\linewidth}}{
        \cellcolor[rgb]{0.9,0.9,0.9}{
            \makecell[{{p{\linewidth}}}]{
                \texttt{\tiny{[GM$|$GM]}}
                {'word': 'buceador', 'assignment': 'team'}
            }
        }
    }
    & & \\ \\

    \theutterance \stepcounter{utterance}  
    & & \multicolumn{4}{p{0.6\linewidth}}{
        \cellcolor[rgb]{0.9,0.9,0.9}{
            \makecell[{{p{\linewidth}}}]{
                \texttt{\tiny{[GM$\rangle$P2]}}
                \texttt{Proporciona una lista separada por comas de hasta 2 palabras de la siguiente lista que mejor se relacionen o estén más estrechamente asociadas con la palabra 'mar'. Comienza siempre tu lista de suposición(es) con ADIVINANZA: y no incluyas ningún otro texto en tu respuesta.} \\
\\ 
\texttt{científico, tokio, corbata, personal, barco, gancho, buceador, capa, diente, gorra, célula, puerto, tono, dragón, soldado, interruptor, escuela, espalda, columpio, china, lugar, tubo, herradura, maestro, enterrador} \\
            }
        }
    }
    & \\ \\

    \theutterance \stepcounter{utterance}  
    & & & \multicolumn{4}{p{0.6\linewidth}}{
        \cellcolor[rgb]{0.9,0.9,0.9}{
            \makecell[{{p{\linewidth}}}]{
                \texttt{\tiny{[GM$\langle$P2]}}
                \texttt{ADIVINANZA: barco, gancho, buceador, puerto} \\
            }
        }
    }
     \\ \\

    \theutterance \stepcounter{utterance}  
    & & & \multicolumn{2}{p{0.3\linewidth}}{
        \cellcolor[rgb]{0.9,0.9,0.9}{
            \makecell[{{p{\linewidth}}}]{
                \texttt{\tiny{[GM$|$GM]}}
                {'player': 'Guesser', 'type': 'wrong number of guesses', 'utterance': 'barco, gancho, buceador, puerto', 'message': 'You made too many guesses (4). You are only allowed to make 2 guesses!', 'guesses': ['barco', 'gancho', 'buceador', 'puerto'], 'number_of_allowed_guesses': 2}
            }
        }
    }
    & & \\ \\

\end{supertabular}
}

\end{document}
